\section{Future Work}
\label{sec:future_work}

Several promising directions remain:

\subsubsection{Feature-Based Models}
Incorporate the 86 predictor features using:
\begin{itemize}
    \item Gradient-boosted trees (LightGBM, XGBoost) with engineered lag features.
    \item Transformer-based architectures (e.g., Temporal Fusion Transformer) that can attend to both temporal and cross-sectional patterns.
    \item Multi-input recurrent networks that process feature vectors alongside the target sequence.
\end{itemize}

\subsubsection{Weighted Loss Functions}
Implement competition-aligned training by:
\begin{itemize}
    \item Using weighted MSE during deep model training.
    \item Developing differentiable approximations to the weighted skill score for end-to-end optimization.
\end{itemize}

\subsubsection{Ensemble Methods}
Combine models using:
\begin{itemize}
    \item Per-series model selection based on validation performance.
    \item Stacking with a meta-learner trained on validation predictions.
    \item Bayesian model averaging weighted by per-series posterior probabilities.
\end{itemize}

\subsubsection{Global and Panel Models}
Move beyond per-series training:
\begin{itemize}
    \item Train a single model across all series with entity embeddings (similar to N-BEATS or DeepAR).
    \item Use hierarchical models that share parameters across related series while allowing series-specific adjustments.
\end{itemize}

\subsubsection{Advanced Architectures}
Explore state-of-the-art time series models:
\begin{itemize}
    \item \textbf{Transformer models:} Self-attention can capture long-range dependencies without recurrence.
    \item \textbf{N-BEATS:} Interpretable deep learning architecture designed for time series forecasting.
    \item \textbf{Neural basis expansion:} Models that decompose forecasts into interpretable trend and seasonality components.
\end{itemize}

\subsection{Closing Remarks}

This project demonstrates that rigorous financial time series forecasting requires a multi-faceted approach: thorough exploratory analysis, principled baseline establishment, and progressive model sophistication. While no single approach achieves dominant performance across the heterogeneous panel, the systematic methodology provides a solid foundation for further improvement through ensemble methods and feature-based modeling.

The competition's design---with anonymized data, weighted evaluation, and strict temporal splitting---mirrors the challenges of real-world quantitative finance, where model robustness and out-of-sample generalization are paramount.

